\section{Hub Design}\label{sec:hub:design}

The hub connects $k$ models through one shared space, with one source adapter and one projector per model (Figure~\ref{fig:hub:arch}).

\begin{figure}[ht]
\centering
\includegraphics[width=\textwidth]{ch3/hub_architecture}
\caption{The translation hub maps each model's hidden states into one shared space, and every frozen model decodes from it.
Model-specific adapters and a shared backbone produce the hub output; model-specific projectors map it into the input space of each frozen decoder.
Solid boxes are trained modules, dashed boxes are frozen models, and the column of cells is the hub output, a vector with no parameters of its own.}
\label{fig:hub:arch}
\end{figure}

\paragraph{Into the hub.} The source adapter of model $a$ maps its hidden state $s_t$ to a common width $d$; a shared backbone then produces the hub output $m_t$ of width $d_h$:
\begin{equation}
m_t = W_h\, R_0\big(R_a(W_a s_t + c_a)\big) + c_h,
\qquad
R(z) = z + V_2\,\mathrm{SiLU}\big(V_1\,\mathrm{LN}(z) + c_1\big) + c_2 .
\label{eq:hub:map}
\end{equation}
Here $W_a$ is the linear layer of model $a$'s adapter, $R_a$ its residual block, $R_0$ the residual block of the shared backbone, and $W_h$ the shared linear layer from width $d$ to width $d_h$.
Each residual block $R$ applies layer normalization (LN), a linear layer $V_1$ that widens to $2d$, a SiLU activation, and a linear layer $V_2$ that returns to $d$; $c_1$ and $c_2$ are the biases of the two layers.
Every block holds its own $V_1$, $V_2$, $c_1$, and $c_2$.
The layout, model-specific adapters followed by a shared backbone, follows Jha et al.~\cite{jha2025harnessing}.
The shared backbone gives every model the same last part of the map, so the hub adds only one adapter per model.

\paragraph{Out of the hub.} A model-specific affine projector maps each hub output $m_t$ to the input width of a frozen language model, the decoder.
The projected vectors form a soft prompt~\cite{lester2021power}, and the decoder generates the reconstructed prompt after its start token.

\paragraph{Directions.} A direction pairs a source model with a decoder: model $a$ $\to$ hub $\to$ model $b$.
The $k$ models define $k^2$ directions: $k(k-1)$ cross-model directions and $k$ self-reconstruction directions.
In a cross-model direction the decoder reconstructs the prompt from another model's hidden states, so the cross-model directions measure translation between models.
Only the adapters, the backbone, and the projectors are trainable; the language models remain frozen as both sources and decoders.

\paragraph{Token-wise mapping.} The hub maps each token independently.
Equation~\ref{eq:hub:map} has no positional input and no attention across tokens, so each hub output depends on one source hidden state, although that state may already encode prompt context.

